%% ---------------------------------------------------------------------------
%% When Does an LLM Beat a Bag of Words at Triaging Static-Analysis Alerts?
%% Computers & Security (Elsevier) -- elsarticle, journal two-column layout
%%
%% DO NOT EDIT main.tex. Edit THIS file, then run:
%%     python3 build.py ../results
%%     pdflatex main && bibtex main && pdflatex main && pdflatex main
%% ---------------------------------------------------------------------------
\documentclass[preprint,12pt]{elsarticle}

\usepackage[htt]{hyphenat}
%% OT1 thay vi T1: Computer Modern ma hoa T1 (ec) chi co ban bitmap
%% Metafont tren he nay, sinh ra chu Type 3 nhoe trong PDF. OT1 da co
%% san ban Type 1 vector. T5 (tieng Viet) dung VNR, cung la Type 1.
\usepackage[T5,OT1]{fontenc}
\usepackage{booktabs}
\usepackage{graphicx}
\usepackage{amsmath}
\usepackage{amssymb}

%% Ngat dong: ten file va dinh danh dai trong \texttt la mot token lien,
%% khong co cho ngat nen bi day ra ngoai le cot.
\sloppy
\emergencystretch=3em
\tolerance=2000

\usepackage{array}
\usepackage{url}

\biboptions{authoryear}
\journal{IEEE Transactions on Software Engineering}
\renewcommand{\figurename}{Fig.}

\newcommand{\vn}[1]{{\fontencoding{T5}\selectfont #1}}

\newcommand{\pendingfig}[2]{%
  \IfFileExists{figures/#1.pdf}%
    {\includegraphics[width=\linewidth]{figures/#1.pdf}}%
    {\fbox{\parbox[c][#2][c]{\dimexpr\linewidth-2\fboxsep-2\fboxrule\relax}%
       {\centering\sffamily\small
        [\,\texttt{\detokenize{figures/#1.pdf}}\,]\par\vspace{3pt}
        \footnotesize hand-drawn diagram --- drop the PDF in
        \texttt{figures/} under this exact name\par}}}}

%%MACROS%%

\usepackage{orcidlink}
\begin{document}

\begin{frontmatter}

\title{We Fabricated a Negative Class Twice: A Failed Repair of a
Static-Analysis Alert-Triage Benchmark}

\author[thinh]{Vinh Thinh Le\,\orcidlink{0000-0001-5951-096X}}
\ead{thinhlv@hcmute.edu.vn}
\author[rmit]{Son X. Ha\,\orcidlink{0000-0002-5336-0566}}
\ead{ha.son@rmit.edu.vn}
\author[tdtu1,tdtu2]{Nguyen Ngoc Phien\,\orcidlink{0000-0001-9152-6208}\corref{cor}}
\ead{nguyenngocphien@tdtu.edu.vn}
\author[fpt]{Tung Q. Nguyen\,\orcidlink{0009-0003-5046-5535}}
\ead{NguyenQuangTung0902@gmail.com}

\cortext[cor]{Corresponding author.}

%% Keep affiliation values plain ASCII on one line: elsarticle parses them as
%% key/value options and a font-switching macro inside a value breaks it.
\affiliation[thinh]{organization={Faculty of Information Technology, Ho Chi Minh City University of Technology and Engineering (HCMUTE)}, addressline={1 Vo Van Ngan Street}, city={Ho Chi Minh City}, postcode={70000}, country={Vietnam}}
\affiliation[rmit]{organization={The Business School, RMIT University}, city={Ho Chi Minh City}, country={Vietnam}}
\affiliation[fpt]{organization={FPT University, Ho Chi Minh City Campus}, city={Ho Chi Minh City}, country={Vietnam}}
\affiliation[tdtu1]{organization={Center for Applied Information Technology \& Faculty of Information Technology, Ton Duc Thang University}, city={Ho Chi Minh City}, country={Vietnam}}
\affiliation[tdtu2]{organization={Faculty of Information Technology, Ton Duc Thang University}, city={Ho Chi Minh City}, country={Vietnam}}

\begin{abstract}
Static analysis produces more alerts than developers can read, and a growing
line of work proposes automated triage as the filter. This paper reports a
negative result about how that work is measured, and a corrected measurement.

The false-positive half of our own earlier evaluation set was not real. Its
$120$ records are $10$ template messages, one per rule identifier, each cloned
exactly $12$ times. A bag-of-words model separates ten fixed strings from a real
positive class almost perfectly, so every comparison drawn on that set,
including ours, is void.

We then attempted a repair, and the repair failed in the same way. Believing we
were scanning third-party code certified as secure, we obtained \RNfp{} alerts
and re-measured four triagers on them. Auditing the provenance afterwards showed
that \CFpSecure{} of those \RNfp{} alerts came from files generated locally for
the purpose: ten directories of twelve near-identical clones, the same
$10 \times 12$ structure, the same \texttt{fp\_N.py} naming and the same ten
distinct messages as the class we had just condemned. The second negative class
is a reconstruction of the first.

This paper reports that failure rather than the result we thought we had. We give
the measurements as they stand, label them as measurements on a generated corpus,
and identify the three conditions that made the same mistake easy to repeat: a
provenance claim that no artefact records, a directory that looks like part of an
upstream benchmark once it sits inside the checkout, and a diversity statistic
that flags the fabricated class and the replacement equally
(\CMultFp{} against $12.00$ alerts per message template). The LLM arm is
withdrawn, not re-run, and no triager comparison in this paper should be quoted
as a benchmark result.

\end{abstract}

\begin{keyword}
Static application security testing \sep CodeQL \sep Alert triage \sep
False positives \sep Benchmark validity \sep Class imbalance \sep
Developer burden

\end{keyword}

\end{frontmatter}

%% ===========================================================================
\section{Introduction}
\label{sec:intro}

Static analysers find real bugs and bury them. The reason developers abandon
static analysis is rarely that it misses vulnerabilities; it is that the alert
stream is dominated by findings that turn out not to matter
\citep{johnson2013why}. Anything that reliably removes false positives without
removing true ones is therefore worth a great deal, and a growing literature
proposes learned or language-model triage as that filter.

Such a filter is only as meaningful as the benchmark it is measured on, and a
triage benchmark has an unusual failure mode: the negative class is the
expensive half. True positives can be taken from any vulnerability corpus, but
genuine false positives have to be produced by running the analyser on code that
is actually secure, and it is tempting to synthesise them instead.

We did exactly that, and this paper reports the consequence. The negative class
of our earlier evaluation set consists of $10$ template messages, each cloned
$12$ times into files named \texttt{CWE-NNN/fp\_N.py}. A classifier that
separates that from a real positive class is separating ten fixed strings from a
population, which is not a triage result. Every comparison we previously drew on
that set is withdrawn.

The rest of the paper was written as the repair. It is now the report of a
repair that failed. We believed we had replaced the generated negative class by
scanning \texttt{SecurityEval}'s secure directories; \texttt{SecurityEval} has no
such directories, and the files we scanned were generated locally two days after
the repository was cloned, by the same pattern as the class they were meant to
replace. Section~\ref{sec:secondfailure} documents that, and
Section~\ref{sec:whytwice} asks why it was so easy to do twice.

We report the measurements we made, because they are real measurements of a real
analyser on real files, and they are what a reader needs in order to judge the
failure. We do not report them as a benchmark result. They are not one.

\paragraph{Findings}
\begin{enumerate}
\item \textbf{The published negative class was generated, not scanned.} Ten
distinct messages across $120$ records, $12$ clones each, in files named
\texttt{CWE-NNN/fp\_N.py} (Section~\ref{sec:fabricated}).
\item \textbf{So was the replacement.} \texttt{SecurityEval} ships no secure
directories. The \CFpSecure{} alerts we took for third-party false positives came
from locally generated files with the same $10 \times 12$ structure and the same
\texttt{fp\_N.py} naming as the class they replaced; only \CFpPatched{} of the
\RNfp{} came from hand-written code (Section~\ref{sec:secondfailure}).
\item \textbf{A lexical classifier clears a hand-written heuristic; the
heuristic does not clear chance.} $\mathrm{MCC} = \RTfidfBalMcc$ against
$\RHeurBalMcc$ with an interval containing zero, paired
$p = \RHeurVsTfidf$ (Section~\ref{sec:realresults}).
\item \textbf{The rule identifier contributes nothing.} Trained on the message
alone the classifier changes at most one decision on either set
($p = \RMsgSame$), which answers the shortcut objection by measurement rather
than by argument (Section~\ref{sec:realresults}).
\item \textbf{The burden metric decides the comparison before any triager
runs.} Charging a missed true positive fifty reviews makes accept-all optimal
at \RAllBalBur{} against \RTfidfBalBur{} for the best classifier
(Section~\ref{sec:realresults}).
\item \textbf{The repaired benchmark still contains two shortcuts, one of
which bounds our own headline number.} The CWE directory settles
\DTpOnOnlyPct\% of the positive class and we show it never reaches the
classifier; the negative class carries only \CMsgFp{} message templates for
\RNfp{} alerts, a multiplicity of \CMultFp{} against \CMultTp{} on the positive
side, which under \(5\)-fold cross-validation puts instances of a test
template in the training folds (Sections~\ref{sec:cweleak}
and~\ref{sec:diversity}). We therefore read
$\mathrm{MCC} = \RTfidfBalMcc$ as a floor any triager must clear and not as a
generalisation estimate.
\item \textbf{The heuristic's failure is attributable clause by clause.} Its
path-keyword filter, which would do most of the work in a real repository,
fires \CPathClauseTp{} times here because the corpus names files by CWE index;
the remaining hedge-word clause rejects \CHedgeTp{} true and \CHedgeFp{} false
positives, which is the whole of the discrimination it achieves
(Section~\ref{sec:heurwhy}).
\end{enumerate}


%% ===========================================================================
\section{Background and Related Work}
\label{sec:related}

\subsection{Why static-analysis alerts need triage at all}

Sound-by-construction static analysis is imprecise by design: to avoid missing
real defects it reports on paths that no execution takes. At industrial scale
the consequence is a queue of alerts that is mostly noise. \citet{bessey2010billion}
report from a decade of commercial deployment that developer tolerance for false
positives collapses well before the rates analyses actually achieve;
\citet{sadowski2018google} describe the organisational machinery Google needed to
keep analysis results credible; and \citet{ayewah2008findbugs} document the same
tension for FindBugs. Interview studies confirm it is the dominant adoption
barrier: \citet{johnson2013why} find false positives and poor presentation are
why developers abandon these tools, and \citet{christakis2016want} find the same
priorities among Microsoft engineers.

The response has been a long line of \emph{actionable alert identification}
techniques, surveyed by \citet{heckman2011slr} and \citet{muske2016survey}.
Statistical ranking of alerts by how often similar ones were acted on goes back
at least to Z-ranking~\citep{kremenek2003zranking}. Modern query-based engines
such as CodeQL~\citep{avgustinov2016ql} produce structured, machine-readable
alerts, which makes automated triage a well-posed classification problem: given
an alert and its code context, predict whether a human would act on it.

\subsection{Learned vulnerability detection and its measurement problems}

A parallel literature learns to detect vulnerabilities directly from code, with
graph and sequence models~\citep{zhou2019devign, li2018vuldeepecker} trained on
curated corpora~\citep{fan2020bigvul}. Two replication studies temper the
reported gains. \citet{chakraborty2022deep} show that models reported as strong
on curated benchmarks degrade sharply on realistic distributions, and
\citet{steenhoek2023empirical} find substantial run-to-run variance and
sensitivity to data preparation. \citet{croft2023dataquality} trace much of the
inconsistency to label quality in the datasets themselves.

The methodological lesson --- that a headline number is only as good as the
data it was measured on: is the one this paper runs into directly, first in
the construction of the negative class and then in the prevalence at which the
comparison is made.

\subsection{Large language models on security-relevant code}

Code-trained language models~\citep{chen2021codex} have been evaluated for
vulnerability repair~\citep{pearce2023repair}, for secure-coding
behaviour~\citep{bhatt2023cyberseceval, siddiq2022securityeval}, and for
detection. The detection results are mixed in an instructive way.
\citet{ullah2024llmcannot} evaluate a broad set of models under a framework
designed to separate correct answers from correct reasoning, and conclude that
current models do not reliably identify or explain vulnerabilities;
\citet{khare2023understanding} find that performance depends heavily on prompting
strategy and on the benchmark's construction. Neither study contrasts such a model
against a cheap lexical baseline on the same alerts under two prevalence
regimes. That comparison motivated this work; the baseline it requires is what
this paper establishes.

\subsection{Bag-of-words baselines}

A TF--IDF representation~\citep{salton1988tfidf} with a linear
classifier~\citep{joachims1998svm} remains a strong baseline for text
classification, and \citet{wang2012baselines} showed that simple
bigram models beat many more elaborate systems when the task has strong lexical
cues. Static-analysis alerts have exactly such cues: rule identifiers, sink
names, and message templates recur verbatim. Any claim that a language model
adds value here has to clear that baseline, on the distribution that matters.

\subsection{Evaluating classifiers under class imbalance}

The prevalence of true positives in a real alert queue is low, and the
consequences for evaluation are well understood but routinely ignored.
\citet{davis2006prroc} and \citet{saito2015prplot} show that ROC curves are
insensitive to prevalence while precision-recall curves are not.
\citet{he2009imbalanced} survey the broader
problem and \citet{chawla2002smote} give the canonical resampling response.
Single-number summaries that behave sensibly under imbalance ---
MCC~\citep{matthews1975, chicco2020mcc}: are preferred here to accuracy,
which a constant predictor can maximise.

Confidence in a prediction is a separate axis from correctness.
\citet{guo2017calibration} showed modern neural classifiers are systematically
overconfident; expected calibration error~\citep{naeini2015ece} and the Brier
score~\citep{brier1950} quantify it. That axis matters for a triager consumed
as a priority ordering, and we do not measure it here: the arm it applies to is
withdrawn.

\subsection{Statistical comparison of classifiers}

Comparing two classifiers on the same test set is a paired problem, so we use
McNemar's test~\citep{mcnemar1947} in its exact binomial form rather than the
$\chi^2$ approximation with a continuity correction~\citep{edwards1948note},
since the discordant counts here are small; \citet{dietterich1998tests} discusses
why the paired form is the appropriate one at this sample size. Proportions
carry Wilson intervals~\citep{wilson1927interval}.

\subsection{Benchmarks as infrastructure, and the cost of getting one wrong}

The failure this paper reports is not a modelling failure, and the literature
that describes it best is the one on machine-learning systems rather than on
machine-learning methods. Sculley et al.~\citep{sculley2015debt} argue that the
code around a model accumulates debt faster than the model itself, and that data
dependencies are the least visible form of it: a dataset generated once by a
script nobody re-reads becomes a fixed input that every later result silently
inherits. That is exactly the shape of our error. The generator that produced
the synthetic negative class was written to unblock an early experiment, ran
correctly, and was never revisited; every comparison built on top of it was
methodologically sound and substantively empty.

The remedy we would propose, and the one we would have wanted applied to us,
is narrower than it first appears. The obvious candidate: flag any class
whose members collapse onto a few distinct strings --- does not work, and the
numbers here are what rule it out. The fabricated negative class held \(120\)
records over \(10\) distinct messages, a multiplicity of \(12.00\). The
scanned replacement holds \RNfp{} records over \CMsgFp{}, a multiplicity of
\CMultFp{}. On that statistic the repaired benchmark is indistinguishable from
the one it replaces, and a threshold tuned to reject the first would reject the
second. Low textual diversity is a property of running one query suite over a
directory of near-identical files, and that is what honest scanning of a
synthetic corpus produces.

We wrote the paragraph above believing that provenance was what separated the
two sets: one emitted by a generator, the other by an analyser running on
third-party code certified as secure. Section~\ref{sec:secondfailure} shows both
sets were emitted by a generator, so the two are separated by nothing at all. We
leave the reasoning in place because the conclusion it reaches survives its own
premise being wrong, and reaches further than we intended.

Provenance --- \emph{scanned}, \emph{annotated} or \emph{generated}, recorded per
class as a property of the benchmark rather than of whichever paper introduces it: is the field we want, and it is a one-word field. But recording it is not
enough if the person filling it in is the person who was fooled. Our pipeline did
carry a provenance field; every alert has a \texttt{source} naming its directory,
faithfully propagated through every stage. The label was correct and the belief
behind it was not. What is needed is a provenance claim an artefact can falsify:
an upstream commit hash, and the assertion that every scanned path is tracked at
it. Section~\ref{sec:whytwice} develops this.

\subsection{Benchmark contamination}

A model that has seen a benchmark during pre-training can appear to reason about
it. \citet{magar2022contamination} separate memorisation from exploitation,
\citet{sainz2023contamination} argue contamination must be measured per
benchmark rather than assumed absent, and \citet{golchin2024timetravel} give a
practical probing method. Our corpus is drawn from public code, so any
pre-trained triager evaluated on it will need such a probe. None of the triagers
compared here is pre-trained, so the question does not arise for these results,
and the probe we previously reported is withdrawn with the rest of the
language-model arm.

%% ===========================================================================
\section{Method}
\label{sec:method}

\begin{figure}[t]
\centering
\pendingfig{fig_pipeline}{46mm}
\caption{
Benchmark construction. \textbf{(A)} True positives from the vulnerable corpus,
false positives from the secure corpus. \textbf{(B)} The two evaluation sets,
differing only in prevalence.
}
\label{fig:pipeline}
\end{figure}

\subsection{Corpus and alert generation}
We run CodeQL's Python security queries over the SecurityEval corpus
\citep{siddiq2022securityeval}, which ships intentionally vulnerable
implementations and no secure or patched counterparts. The vulnerable corpus
yields \RNtp\ true-positive alerts spanning \RNrulesTp\ distinct rules; the most
frequent are listed in Table~\ref{tab:alerts}. The negative class is not
upstream: we generated the secure files inside our own checkout, and they yield
\RNfp\ alerts that are false positives by construction, spanning \RNrulesFp\
rules with \RNmsgFp\ distinct messages, \RNrulesShared\ of which also raise true
positives.
Figure~\ref{fig:pipeline} summarises the construction, and
Section~\ref{sec:fabricated} explains why the negative class must be obtained
this way and what happens when it is not.

\begin{table}[t]
\caption{Ten most frequent rules among the \RNtp{} true-positive alerts.}
\label{tab:alerts}
\centering
\footnotesize
\begin{tabular}{lr}
\toprule
Rule & Alerts \\
\midrule
%%ROWS:tab_alerts%%
\bottomrule
\end{tabular}
\end{table}

\subsection{Two benchmarks}
Prevalence is the only variable we manipulate. The balanced set takes every true
positive and an equal number of false positives, \RNbal\ alerts. The imbalanced
set takes every false positive and subsamples true positives to $1{:}\RImbRatio$,
\RNimb\ alerts, closer to what an unfiltered queue looks like. Each set exhausts
whichever class limits it rather than being capped at a round number.

\subsection{Triagers}
\emph{Accept-all} forwards every alert. It is the no-triage control and fixes
the upper bound on developer burden.

\emph{Heuristic regex triage} applies hand-written patterns per rule family, the
kind of filter a team writes after a few weeks of alert fatigue.

\emph{TF-IDF + Random Forest} vectorises \texttt{rule\_id} and the alert message
and classifies with a Random Forest under five-fold cross-validation. This is
the baseline any learned or language-model triager has to clear, because it uses
no semantics beyond token frequency.

\emph{TF-IDF + Random Forest, message only} is the same model with
\texttt{rule\_id} removed from the input. It exists to test whether the
classifier is recovering the class from the rule identifier rather than from the
alert text.

\subsection{Metrics}
\label{sec:metrics}
Alongside accuracy and the confusion matrix we report the Matthews correlation
coefficient \citep{matthews1975}, which is the appropriate summary under class
imbalance \citep{chicco2020mcc} and which a constant predictor cannot maximise.

We also report a \emph{developer burden}
$B_{\mathrm{dev}} = \mathrm{FP} + 50\,\mathrm{FN}$, a cost in review-units
combining the alerts a developer must still read with a penalty for missed true
positives, so that a triager cannot look good by rejecting everything. The
multiplier is inherited from the surrounding literature;
Section~\ref{sec:realresults} shows what it does to the comparison.

\subsection{Statistics}
Paired comparisons between triagers on the same alerts use an exact two-sided
McNemar test \citep{mcnemar1947}, computed from the binomial distribution of the
discordant pairs rather than from a $\chi^2$ approximation, since the discordant
counts here are small. Intervals on MCC are percentile bootstrap intervals over
\RBoot\ resamples of the evaluation set. No multiplicity correction is applied;
the comparisons we interpret are stated in advance and reported with their raw
$p$-values.

%% ===========================================================================
\section{The Negative Class Was Not Real}
\label{sec:fabricated}

Before any triager can be compared, the benchmark it is compared on has to
exist. Ours did not, and we report that first because it invalidates the
comparison the first version of this paper made.

The false-positive half of the original evaluation set was read from a file of
$120$ alert records. Those records carry $10$ distinct alert messages, one per
rule identifier, each repeated exactly $12$ times, in $120$ files named
\texttt{CWE-NNN/fp\_N.py} across $10$ CWE directories. The regularity is
mechanical: every rule contributes the same count, and every alert of a rule
carries character-for-character the same message. They were generated, not
scanned. A classifier separating that negative class from a genuine positive
class is separating ten fixed strings from a population of real ones, which any
bag-of-words model does almost perfectly and which tells an auditor nothing.

We set out to rebuild the negative class by scanning code known to be secure.
We believed \texttt{SecurityEval} shipped two directories of deliberately safe
implementations, \texttt{Testcases\_Secure\_FP} and
\texttt{Testcases\_Secure\_Patched}, and that every CodeQL alert raised on them
was therefore a false positive by construction. The scan produced \RNfp{} alerts
across \RNrulesFp{} rules with \RNmsgFp{} distinct messages, against \RNtp{}
true positives across \RNrulesTp{} rules from the vulnerable corpus.

That belief was false, and Section~\ref{sec:secondfailure} sets out how we
established it. \texttt{SecurityEval} publishes six directories ---
\texttt{Testcases\_Prompt}, \texttt{Testcases\_Insecure\_Code},
\texttt{Testcases\_Copilot}, \texttt{Testcases\_InCoder}, \texttt{Databases} and
\texttt{Result}: and no secure implementations at all. The two directories we
scanned were created inside our own checkout after the clone.

\paragraph{The rule identifier is not a shortcut.}
Within the inventory as measured, a classifier cannot win by memorising which
rule identifiers belong to which class: \RNrulesShared{} of the \RNrulesFp{}
negative-class rules also appear among the true positives. We test this by
running the same classifier twice, once on \texttt{rule\_id} plus message and
once on the message alone, and Section~\ref{sec:realresults} reports that the two
are statistically indistinguishable. We record the check because it is the one
objection we did anticipate: and because anticipating it, and answering it
with a measurement, is exactly what gave us confidence in an inventory whose
provenance we had never checked.

%% ===========================================================================
\section{The Replacement Was Generated Too}
\label{sec:secondfailure}

Everything up to here was written believing the repair had worked. This section
is the audit that showed it had not, and it is the paper's actual contribution.

\subsection{The upstream benchmark has no secure directories}

\texttt{SecurityEval}~\citep{siddiq2022securityeval} is a dataset of Python
code-generation prompts paired with insecure completions. Its published
repository contains \texttt{Testcases\_Prompt},
\texttt{Testcases\_Insecure\_Code}, \texttt{Testcases\_Copilot},
\texttt{Testcases\_InCoder}, \texttt{Databases} and \texttt{Result}. There is no
\texttt{Testcases\_Secure\_FP} and no \texttt{Testcases\_Secure\_Patched}, and
the benchmark ships no secure or patched implementations of any kind. The
property our repair depended on does not exist, and the authors we credited with
certifying it certify nothing, because they never wrote these files.

\subsection{The directories were created inside our checkout}

Both directories exist on our disk. Their modification times place them after
the clone: the repository checkout carries one timestamp,
\texttt{Testcases\_Secure\_Patched} is \(2.4\) days later and
\texttt{Testcases\_Secure\_FP} is \(2.6\) days later. They were produced locally,
during this project, and then read back as though they were upstream material.

\subsection{Their contents are clones, with the analyser in mind}

\texttt{Testcases\_Secure\_FP/CWE-605/} holds \texttt{fp\_1.py} through
\texttt{fp\_12.py}. Each is \(547\) or \(548\) bytes and each is byte-identical
to the others except for the index embedded in its function name. Each carries a
comment naming the rule it is meant to trip. The files were written to produce
alerts, not to be secure code that happens to attract them.

\subsection{The two negative classes are the same object}

Table~\ref{tab:twoclasses} places the class we condemned beside the class we
replaced it with. Every structural statistic matches.

\begin{table}[t]
\caption{The fabricated negative class and its replacement, compared on
properties that do not depend on any modelling choice. The replacement's
\texttt{secure\_fp} component reproduces the original exactly.}
\label{tab:twoclasses}
\centering
\footnotesize
\setlength{\tabcolsep}{5pt}
\begin{tabular}{@{}lrr@{}}
\toprule
Property & Fabricated class & Replacement (\texttt{secure\_fp}) \\
\midrule
Alerts & 120 & \CFpSecure \\
CWE directories & 10 & 10 \\
Clones per directory & 12 & 12 \\
File naming & \texttt{fp\_N.py} & \texttt{fp\_N.py} \\
Distinct messages & 10 & 10 \\
Distinct rules & 10 & 10 \\
\bottomrule
\end{tabular}
\end{table}

Of the \RNfp{} alerts in the replacement, \CFpSecure{} come from those generated
clones. The remaining \CFpPatched{} come from \texttt{Testcases\_Secure\_Patched},
whose files are hand-written and genuinely varied: and which were also written
locally rather than obtained from a third party. So the honest count of the
negative class is \CFpPatched{} alerts on code we wrote ourselves, and
\CFpSecure{} on code we generated for the purpose.

Every triager comparison in this paper is therefore a comparison on a generated
negative class. We report the numbers because they are correct measurements of
what was actually scanned, and because withholding them would leave the failure
undocumented. They are not a benchmark result and must not be cited as one.

%% ===========================================================================
\section{Why It Happened Twice}
\label{sec:whytwice}

The interesting question is not that a benchmark was fabricated. It is that the
team that discovered the fabrication, wrote a paper about it, and built a
correction, produced the same artefact again. Three conditions made that
possible, and all three are general.

\paragraph{Provenance is asserted in prose and recorded nowhere.}
The scan script states in its own docstring that
``SecurityEval co san hai thu muc code AN TOAN'' and that any alert on them is a
false positive because the benchmark's authors vouch for the code. Nothing
downstream checks that. The alert records carry a \texttt{source} field naming
the directory, and that field was faithfully propagated through every later
stage, which made the provenance feel documented when only its label was. A
provenance claim that no artefact can falsify is a comment, not a record.

\paragraph{A generated directory inside a checkout looks like part of it.}
Once \texttt{Testcases\_Secure\_FP} sat next to
\texttt{Testcases\_Insecure\_Code} under \texttt{repos/SecurityEval/}, every path
in every log read as upstream material. Nothing in the pipeline distinguishes a
directory that arrived with \texttt{git clone} from one written afterwards, and
neither did we. Checking costs one command --- \texttt{git status} would have
listed both directories as untracked from the moment they appeared.

\paragraph{The diversity check we proposed does not separate the two classes.}
Section~\ref{sec:diversity} reports the negative class as \CMsgFp{} message
templates over \RNfp{} alerts, a multiplicity of \CMultFp{}. The fabricated class
was \(120\) records over \(10\) templates, \(12.00\). We had written that ratio
up as the honest consequence of scanning a synthetic corpus. A threshold that
rejects the first rejects the second, and one that accepts the second accepts the
first. The statistic we offered as a safeguard would have passed the fabricated
class, and our own low reading of it was the clue we explained away.

\paragraph{What we would require instead.}
Provenance has to be a property of the artefact, verifiable without trusting the
text. Three checks would each have caught this on their own: record the upstream
commit hash and assert that every scanned path is tracked at that commit; store,
per class, whether its files came from a clone, from an annotator, or from a
generator, and refuse to build an evaluation set from a class whose provenance is
unrecorded; and treat any directory whose file sizes are near-constant and whose
names are an index sequence as generated until shown otherwise. None of these
requires domain knowledge. All three are cheaper than the audit that eventually
caught us, which began only because an outside reviewer asked where the secure
code came from.

%% ===========================================================================
\section{What the Reconstructed Inventory Contains}
\label{sec:inventory}

Before any classifier is run it is worth stating precisely what the two classes
are made of, because two properties of the inventory bound what any result on it
can mean. Both properties are read directly off the two alert files and involve
no modelling choice.

\subsection{Rule identifiers overlap heavily between the classes}

The \CNall{} alerts carry \CNtpRules{} distinct rule identifiers on the
true-positive side and \CNfpRules{} on the false-positive side, and
\CNshared{} identifiers occur on both. That overlap is not marginal.
\CFpOnShared{} of the \RNfp{} false positives (\CFpOnSharedPct\%) and
\CTpOnShared{} of the \RNtp{} true positives (\CTpOnSharedPct\%) sit on a rule
that also produces alerts of the other class, so for roughly half the positives
and nine tenths of the negatives the rule identifier carries no class
information at all.

Table~\ref{tab:ruleoverlap} lists the shared identifiers with their counts in
each class. The imbalance within a rule varies by almost two orders of
magnitude: \texttt{py/reflective-xss} is overwhelmingly a true positive in this
corpus (7 against 1), while \texttt{py/bind-socket-all-network-interfaces},
\texttt{py/incomplete-hostname-regexp} and
\texttt{py/request-without-cert-validation} are overwhelmingly false positives
(1 against 12 each). A triager that learned a per-rule prior would therefore be
right about some rules and badly wrong about others, which is the quantitative
version of the qualitative claim that alert triage cannot be reduced to a rule
allow-list.

\begin{table}[t]
\centering
\caption{Rule identifiers that produce alerts in both classes. \textbf{TP} and
\textbf{FP} are counts in the reconstructed inventory; the last column is their
ratio. Read from \texttt{pilot/alerts.csv} and
\texttt{real\_fp/real\_fp\_alerts.json}.}
\label{tab:ruleoverlap}
\footnotesize
\setlength{\tabcolsep}{4pt}
\begin{tabular}{@{}lrrr@{}}
\toprule
Rule identifier & TP & FP & FP/TP \\
\midrule
%%ROWS:tab_rule_overlap%%
\midrule
Total on shared rules & \CTpOnShared & \CFpOnShared & \\
\bottomrule
\end{tabular}
\end{table}

\subsection{The CWE directory is a near-perfect label, and the classifier does
not see it}
\label{sec:cweleak}

A third shortcut is worth documenting because it is the kind that silently
inflates published numbers, and because in this case we can show it was not
taken.

The corpus is organised one directory per CWE category, so every alert's
relative path begins with a group name. Table~\ref{tab:cwe} gives the
distribution. The positive class spans \DCweTp{} CWE groups and the negative
class only \DCweFp{}, of which \DCweBoth{} carry both classes. That leaves
\DCweOnlyTp{} groups (holding \DTpOnOnly{} of the \RNtp{} true positives,
\DTpOnOnlyPct\% of the positive class) in which no false positive occurs at
all. A feature that revealed the directory would therefore label
\DTpOnOnlyPct\% of the positives with certainty, and one group
(\texttt{CWE-089}) would label its \(12\) alerts as negatives with equal
certainty. A model given the path would score very well and would have learned
nothing about alert triage.

We can rule this out by inspection rather than by argument, because the feature
construction is explicit: both TF-IDF variants are built from
\texttt{rule\_id} concatenated with the alert message, or from the message
alone. The relative path is carried in the inventory, used for the heuristic's
path-keyword clause (Section~\ref{sec:heurwhy}), and never given to the
learned model. That is why we report the feature text verbatim rather than
describing it: a reader checking whether our numbers are a directory artefact
needs to see the exact string that was vectorised, and no summary of it will
do.

The general point is that a corpus organised by vulnerability class hands every
model trained on it a label-shaped feature, and the organisation is usually
inherited from the benchmark rather than chosen by the person running the
experiment. Whether it leaks depends entirely on which fields reach the
vectoriser, which is a one-line detail that papers routinely omit.

\begin{table}[t]
\centering
\caption{Alerts per CWE group. The \DCweOnlyTp{} groups summarised in the last
row contain \DTpOnOnly{} true positives and no false positive, so the directory
name alone would settle \DTpOnOnlyPct\% of the positive class.}
\label{tab:cwe}
\footnotesize
\setlength{\tabcolsep}{6pt}
\begin{tabular}{@{}lrr@{}}
\toprule
CWE group & TP & FP \\
\midrule
%%ROWS:tab_cwe%%
\midrule
Total & \RNtp & \RNfp \\
\bottomrule
\end{tabular}
\end{table}

\subsection{The negative class is far less textually diverse than the positive}
\label{sec:diversity}

The second property is an asymmetry we did not design for and that materially
affects how the classifier results should be read. The \RNtp{} true positives
carry \CMsgTp{} distinct alert messages, an average of \CMultTp{} alerts per
message template. The \RNfp{} false positives carry only \CMsgFp{} distinct
messages for \RNfp{} alerts, an average of \CMultFp{}. Table~\ref{tab:fptmpl}
gives the full template inventory of the negative class; ten of the twelve
templates account for twelve or more alerts each.

The cause is mechanical. The false positives were harvested by running the same
query suite over a corpus of secure variants organised one directory per CWE,
with several near-identical files per directory
(\CFpSecure{} alerts from the secure corpus and \CFpPatched{} from the patched
variants). Each query emits its message with the same wording every time it
fires, so a directory of twelve structurally similar secure files yields twelve
copies of one template. Nothing about that is wrong as a way to obtain real
false positives --- these are genuine alerts that CodeQL genuinely raised on
code that is genuinely not vulnerable: but it does mean the negative class is
close to twelve repeated strings rather than \RNfp{} independent observations.

We return to the consequence in Section~\ref{sec:limitations}: it is the single
largest threat to the TF-IDF result reported next, and it is the reason we do
not present that result as evidence that lexical triage generalises.

\begin{table}[t]
\centering
\caption{Message templates of the negative class, with the number of alerts each
accounts for. Twelve templates cover all \RNfp{} false positives.}
\label{tab:fptmpl}
\footnotesize
\setlength{\tabcolsep}{4pt}
\begin{tabular}{@{}r>{\raggedright\arraybackslash}p{0.86\linewidth}@{}}
\toprule
$n$ & Alert message template \\
\midrule
%%ROWS:tab_fp_templates%%
\bottomrule
\end{tabular}
\end{table}

%% ===========================================================================
\section{Results on the Rebuilt Benchmark}
\label{sec:realresults}

From the \RNtp{} true and \RNfp{} false positives we build two evaluation sets.
The balanced set uses every true positive and an equal number of false
positives, \RNbal{} alerts. The imbalanced set uses every false positive and
subsamples true positives to a ratio of $1{:}\RImbRatio$, \RNimb{} alerts,
which is closer to what a developer sees. Each set therefore exhausts whichever
class limits it, rather than being capped at a round number.

Four triagers are compared: accepting every alert, a hand-written regex
heuristic, TF-IDF with a Random Forest over \texttt{rule\_id} and message, and
the same model over the message alone. Matthews correlation is reported with a
percentile bootstrap over \RBoot{} resamples, and pairs are compared with an
exact McNemar test on the identical alert set.

\begin{table}[t]
\caption{Balanced set, \RNbal{} alerts. Confusion is TP/FP/TN/FN. Burden is
$\mathrm{FP} + 50\,\mathrm{FN}$, the cost model of
Section~\ref{sec:metrics}.}
\label{tab:realbal}
\centering
\footnotesize
\setlength{\tabcolsep}{3.5pt}
\begin{tabular}{lrlrrlr}
\toprule
Triager & $n$ & TP/FP/TN/FN & Acc. & MCC & 95\% CI & $B_{\mathrm{dev}}$ \\
\midrule
%%ROWS:tab_real_balanced%%
\bottomrule
\end{tabular}
\end{table}

\begin{table}[t]
\caption{Imbalanced set, \RNimb{} alerts at $1{:}\RImbRatio$.}
\label{tab:realimb}
\centering
\footnotesize
\setlength{\tabcolsep}{3.5pt}
\begin{tabular}{lrlrrlr}
\toprule
Triager & $n$ & TP/FP/TN/FN & Acc. & MCC & 95\% CI & $B_{\mathrm{dev}}$ \\
\midrule
%%ROWS:tab_real_imbalanced%%
\bottomrule
\end{tabular}
\end{table}

\begin{figure*}[t]
\centering
\includegraphics[width=\textwidth]{figures/fig_mcc.pdf}
\caption{Matthews correlation with 95\% bootstrap intervals. Only the two
TF-IDF variants separate from zero, on either prevalence.}
\label{fig:mcc}
\end{figure*}

\subsection{The lexical model separates the classes; the heuristic does not}
\label{sec:sep}

On the balanced set the TF-IDF classifier reaches
$\mathrm{MCC} = \RTfidfBalMcc$ with a bootstrap interval of
$[\RTfidfBalLo, \RTfidfBalHi]$, while the regex heuristic reaches
$\RHeurBalMcc$ with an interval of $[\RHeurBalLo, \RHeurBalHi]$: an interval
that contains zero. On the imbalanced set (Table~\ref{tab:realimb}) the gap widens:
$\RTfidfImbMcc$ against $\RHeurImbMcc$, whose interval again contains zero. Paired McNemar
separates the two decisively on both sets
(balanced $p = \RHeurVsTfidf$, imbalanced $p = \RHeurVsTfidfImb$).
Figure~\ref{fig:mcc} shows both regimes with their intervals.

The practical reading is that the alert message carries enough lexical signal
to be worth exploiting, and that a hand-written pattern list does not capture
it. That is a modest claim, but it is the one the data support, and it is the
baseline any learned or LLM-based triager has to clear.

\subsection{Why the heuristic baseline fails, clause by clause}
\label{sec:heurwhy}

The regex heuristic is the baseline the literature's hand-written filters
stand in for, and its near-zero MCC deserves an explanation rather than a
shrug. Because the heuristic is deterministic and reads only the rule
identifier, the message and the file path, we can attribute every one of its
decisions over the whole inventory to the clause that produced it. No sampling
is involved, so unlike the benchmark results this attribution is exact.

Table~\ref{tab:heurclause} gives the attribution. Two facts explain the
outcome. First, the path-keyword clause --- which rejects any alert whose path
contains \texttt{test}, \texttt{example}, \texttt{mock} or \texttt{sample},
and which in a real repository would do most of the filtering: fires
\CPathClauseTp{} times on true positives and \CPathClauseFp{} times on false
positives, that is, never. The corpus names its files by CWE and index rather
than by role, so the clause has nothing to match. Second, the fall-through
default is \emph{accept}, and it catches 45 of \RNtp{} true positives and 90 of
\RNfp{} false positives. The heuristic therefore behaves, on this inventory,
almost exactly like accept-all with a small hedge-word carve-out: the hedge
clause rejects \CHedgeTp{} true positives and \CHedgeFp{} false positives, a
ratio of roughly one to four, which is the only discrimination the whole filter
achieves.

Over the complete inventory --- all \RNtp{} true and all \RNfp{} false
positives, with no subsampling --- the heuristic accepts \CHeurFullTP{} of
\RNtp{} true positives (\CHeurAccTpPct\%) and \CHeurFullFP{} of \RNfp{} false
positives (\CHeurAccFpPct\%), giving
$\mathrm{MCC} = \CHeurFullMcc$ against $\CAllPrec$ precision for accept-all on
the same inventory. That number is not the benchmark result and we do not
report it as one; it is a property of the filter and the corpus, and it is the
mechanism behind the benchmark result.

The general lesson is not that regex filters are useless. It is that a filter
whose discriminating clauses key on repository conventions --- directory names,
test markers, vendor paths --- transfers nothing to a corpus that does not
follow those conventions, and a benchmark built from a synthetic corpus will
therefore understate such a filter's field performance while overstating the
margin any learned model holds over it. Both directions of that bias are worth
naming.

\begin{table}[t]
\centering
\caption{Which clause of the regex heuristic decides each alert, over the whole
inventory. The path-keyword clause never fires on this corpus.}
\label{tab:heurclause}
\footnotesize
\setlength{\tabcolsep}{5pt}
\begin{tabular}{@{}llrr@{}}
\toprule
Clause & Verdict & TP & FP \\
\midrule
%%ROWS:tab_heur_clauses%%
\midrule
Total & & \RNtp & \RNfp \\
\bottomrule
\end{tabular}
\end{table}

\subsection{Dropping the rule identifier changes nothing}

The classifier trained on the message alone reaches $\RMsgBalMcc$ on the
balanced set against $\RTfidfBalMcc$ for the version that also sees
\texttt{rule\_id}, and $\RMsgImbMcc$ against $\RTfidfImbMcc$ on the imbalanced
set. Paired McNemar finds no difference on either
($p = \RMsgSame$ and $p = \RMsgSameImb$; the two disagree on at most one alert).

This settles the shortcut question raised in Section~\ref{sec:fabricated}. The
model is not recovering the class from the rule identifier, because removing
the identifier entirely leaves its decisions unchanged. What it uses is the
message text.

\begin{figure}[t]
\centering
\includegraphics[width=\linewidth]{figures/fig_burden.pdf}
\caption{Recall against developer burden, log scale. The fifty-fold penalty on
misses puts accept-all at the bottom of both curves.}
\label{fig:burden}
\end{figure}

\subsection{The cost model makes accepting everything optimal}

Table~\ref{tab:realbal} contains a result that is awkward for the metric rather
than for any triager. Accepting every alert scores
$\mathrm{MCC} = \RAllBalMcc$ by construction: it never discriminates --- yet
its developer burden is $\RAllBalBur$, the \emph{lowest} of the four, against
$\RTfidfBalBur$ for the classifier with the highest MCC. The same inversion
holds on the imbalanced set, $\RAllImbBur$ against $\RTfidfImbBur$.
Figure~\ref{fig:burden} plots the trade-off: burden falls monotonically with
recall, which is the opposite of what a triage metric should reward.

The cause is the weighting. A burden of
$\mathrm{FP} + 50\,\mathrm{FN}$ penalises a missed vulnerability fifty times
more than a wasted review, so a triager that never rejects anything incurs no
false negatives and therefore no fifty-fold penalty. Under that cost model
accepting every alert is not merely competitive, it is unbeatable, and no
amount of classifier quality can overcome it.

We report this rather than quietly re-weighting, because the same cost model
appears in the surrounding literature and the same inversion will appear
wherever it is used with a large enough false-negative multiplier. A burden
metric is only informative if the multiplier is small enough that discrimination
can pay for itself; otherwise the honest summary is that the metric has already
chosen the answer. Choosing the multiplier is a deployment decision, and it
should be stated and justified rather than inherited.

\subsection{All pairwise comparisons, both prevalences}
\label{sec:allpairs}

Section~\ref{sec:realresults} quotes the two McNemar tests that carry the
argument. For completeness Table~\ref{tab:mcnfull} gives all six pairwise tests
on each set, with the discordant counts $b$ and $c$ that the exact test is
computed from. Reporting the counts matters because a $p$-value on 118 or 148
paired decisions is easy to over-read: the two TF-IDF variants, for instance,
differ on a single alert out of \RNbal{} on the balanced set and one out of
\RNimb{} on the imbalanced set, which is why their $p$-values are exactly
$1$ and why we treat them as one method rather than two.

Three groupings emerge and they are consistent across prevalences. Accept-all
and the heuristic are not separated on the balanced set
($b=7$, $c=15$, $p=0.1338$) and are separated on the imbalanced set only
because the heuristic's hedge clause rejects enough of the now-dominant
negatives to register. Both TF-IDF variants separate from both of them on both
sets, most sharply on the imbalanced set where the discordant counts are
lopsided by an order of magnitude ($15$ against $126$, and $11$ against
$101$). And the two TF-IDF variants do not separate from each other anywhere.
The ordering of the four triagers is therefore stable under a change of
prevalence, even though every individual metric value moves.

\begin{table*}[t]
\centering
\caption{Exact McNemar tests for all pairs of triagers on both evaluation sets.
$b$ is the number of alerts the first method gets right and the second wrong,
$c$ the reverse. Computed by \texttt{fix2\_rebuild\_and\_triage.py} and read
from \texttt{real\_eval/real\_eval\_results.json}.}
\label{tab:mcnfull}
\footnotesize
\setlength{\tabcolsep}{6pt}
\begin{tabular}{@{}lllrrr@{}}
\toprule
Set & Method A & Method B & $b$ & $c$ & $p$ \\
\midrule
%%ROWS:tab_mcnemar_full%%
\bottomrule
\end{tabular}
\end{table*}

\subsection{What is not yet measured}
\label{sec:notyet}

The LLM arm has not been run on the rebuilt benchmark. The earlier version of
this paper reported LLM accuracy, cost, latency, calibration, context-window
ablation and explanation quality on the original evaluation set, whose negative
class we have just shown was fabricated. Those numbers do not transfer and we
have withdrawn them rather than reprint them beside the corrected benchmark.

What that leaves is a comparison among three non-LLM triagers plus an
accept-all control, on a benchmark whose both classes are real. We regard this
as the necessary first step rather than the finished study: the baseline the
LLM must clear is now established and measured, and the LLM measurement against
it is the immediate next experiment.

\section{Discussion}
\label{sec:discussion}

\subsection{Why a synthetic negative class is not a conservative shortcut}
The intuition behind generating false positives is that they are a nuisance
class: the classifier's job is to reject them, so their internal diversity
should not matter much. The intuition is wrong in a specific way. A triager is
scored on how well it separates two populations, and if one of them is a single
point then the measured separation is a property of the generator, not of the
triager. It does not make the benchmark harder or easier in a way that can be
corrected for afterwards; it makes the number meaningless. The only repair is to
obtain the negative class the same way the positive class was obtained, by
running the analyser and reading what it reports.

The repair is also cheap. Corpora built for secure-coding evaluation already
ship safe implementations, and every alert an analyser raises on them is a false
positive by construction, with no annotation effort at all. We see no reason for
a triage benchmark to synthesise its negatives.

\subsection{What the corrected numbers support, and what they do not}
On the rebuilt benchmark the honest statement is narrow. The alert message
carries lexical signal that a Random Forest can exploit,
$\mathrm{MCC} = \RTfidfBalMcc$ on the balanced set against $\RHeurBalMcc$ for a
hand-written pattern list whose interval contains zero. That is a real
difference, paired $p = \RHeurVsTfidf$, and it widens under imbalance. It is
also the full extent of the claim: we have shown that a cheap lexical model
beats a cheap hand-written one, not that either is fit for deployment.

The recall figures make the second half of that sentence concrete. The
classifier recovers $35$ of \RNtp{} true positives on the balanced set and $6$
of $21$ on the imbalanced one. A triager that silently discards two-thirds of
the vulnerabilities it is shown is not a triager a security team can adopt,
whatever its MCC. Bag-of-words separation and usable triage are different
targets, and this benchmark currently measures only the first.

\subsection{A burden metric can decide the comparison before it starts}
The developer-burden result is the one we expect to generalise beyond this
paper. With a fifty-fold penalty on missed true positives, accept-all is
unbeatable by construction: it has no false negatives, so it pays no penalty,
and every discriminating triager pays a penalty that no reduction in review
volume can offset. The metric has chosen its winner before any system runs.

This is not an argument against cost-weighted evaluation, which is the right
frame for triage. It is an argument that the multiplier is a deployment
parameter and has to be stated, justified, and varied. A multiplier small enough
that discrimination can pay for itself yields an informative ranking; one large
enough to dominate yields a tautology. Reporting the ranking at a single
inherited multiplier hides which of the two you are in.

\subsection{Reporting practice}
Three habits would make this literature easier to check. Report how the negative
class was obtained, in the same detail as the positive class, including the
number of distinct alert messages it contains. Report a term-frequency baseline,
which is a few lines of scikit-learn and separates a claim about language
understanding from a claim about lexical classification. Report prevalence
prominently and evaluate at more than one value of it, since it changes the
ranking and not merely the magnitudes.

%% ===========================================================================
\section{Limitations}
\label{sec:limitations}

\emph{The LLM comparison is absent.} The question in our earlier title, whether
a language model beats a bag of words, is not answered here. The arm was not
re-run on the rebuilt benchmark, and the numbers from the fabricated one are
withdrawn rather than carried forward. The baseline that arm must clear is now
established, which is the prerequisite for the measurement, not a substitute
for it.

\emph{Small benchmarks.} \RNbal\ and \RNimb\ alerts. The bootstrap intervals are
correspondingly wide, which is why the heuristic's interval containing zero is
reported as a null result rather than as a small win.

\emph{Template multiplicity is the largest threat to the TF-IDF result.}
This is the limitation we want read before the headline number, because it
bounds what that number means. As Section~\ref{sec:diversity} measures, the
negative class contains \CMsgFp{} distinct message templates across \RNfp{}
alerts, a multiplicity of \CMultFp{}, against \CMultTp{} on the positive side.
The classifier is evaluated by \(5\)-fold stratified cross-validation, so for
any template with twelve or more instances: ten of the \CMsgFp{} templates
qualify --- instances of that same template are present in the training folds
whenever it appears in a test fold. A bag-of-words model can then score a test
alert by recognising a string it has already seen labelled, which is a much
easier problem than the one a deployed triager faces.

We therefore do not claim that $\mathrm{MCC} = \RTfidfBalMcc$ estimates how a
lexical triager would perform on unseen alerts. The defensible reading is
narrower and still useful: on a benchmark whose negative class is real rather
than fabricated, the message text is separable and a hand-written pattern list
does not separate it, so any triager that claims credit must at minimum beat a
lexical model rather than only the regex filter. Turning the number into a
generalisation estimate requires grouped cross-validation that holds whole
templates out, or a second corpus, and we have run neither. We note that the
grouped design would leave at most \CMsgFp{} negative groups, which is too few
to support a stable estimate: so the honest conclusion is that this corpus
cannot answer the generalisation question at all, and a larger negative corpus
is the prerequisite rather than a refinement.

\emph{The negative class comes from one corpus.} Every false positive is a
CodeQL alert on SecurityEval's secure directories, \RNrulesFp{} rules and
\RNmsgFp{} distinct messages, harvested as \CFpSecure{} alerts from the secure
variants and \CFpPatched{} from the patched ones. That is a genuine negative
class, but it is not a diverse one, and a benchmark drawn from a second corpus
would test whether the lexical signal we measure survives a change of source.

\emph{The heuristic is penalised by the corpus, not only by its design.}
Section~\ref{sec:heurwhy} shows its path-keyword clause never fires here,
because the corpus names files by CWE index rather than by role. In a real
repository that clause would carry much of the filter's weight. The heuristic's
\CHeurFullMcc{} on this inventory is therefore a lower bound on what the same
filter would achieve in the field, and correspondingly the margin we report for
the lexical model over it is an upper bound. Neither bound is tight and we do
not know how loose either is.

\emph{The corpus layout is a latent leak for any future arm.} We verify in
Section~\ref{sec:cweleak} that the CWE directory does not reach the two TF-IDF
variants. That verification does not extend to the LLM arm that
Section~\ref{sec:notyet} leaves unrun: a language model prompted with a file
path, a code excerpt, or anything else that carries the directory name would
have access to a feature that settles \DTpOnOnlyPct\% of the positive class.
Whoever runs that arm has to state which fields the prompt contains, and the
result is uninterpretable without it.

\emph{One language, one analyser.} Python, CodeQL. C/C++ with its different
false-positive profile, and other analysers, are untested.

\emph{Labels from limited annotation.} True positives come from manual
inspection of the flagged flows; false positives are labelled by construction
from the secure corpus. We report no inter-annotator agreement statistic because
the annotation was not run redundantly, and we make no claim that would require
one.

\emph{The burden model is a model.} $B_{\mathrm{dev}}$ embeds a chosen penalty
for missed true positives, and Section~\ref{sec:realresults} shows the ranking
is not robust to that choice. The burden columns are not portable to a team with
a different miss cost.

%% ===========================================================================
\section{Conclusion}
\label{sec:conclusion}

The false-positive half of our earlier evaluation set was ten template
messages, each cloned twelve times. Every triage comparison drawn on it is
withdrawn.

We then rebuilt it, and rebuilt the same thing. \CFpSecure{} of the \RNfp{}
alerts in the replacement come from ten directories of twelve near-identical
files generated inside our own checkout, carrying the same
\texttt{fp\_N.py} naming and the same ten messages as the class they replaced.
The remaining \CFpPatched{} come from hand-written files, also ours. No part of
the negative class came from a third party, and the upstream benchmark we
credited ships no secure code at all.

On that inventory a TF-IDF Random Forest reaches
$\mathrm{MCC} = \RTfidfBalMcc$ where a regex heuristic reaches $\RHeurBalMcc$
with an interval containing zero, and under the burden metric the field commonly
uses, accepting every alert beats every triager. Those are correct measurements
of what was scanned and they are not benchmark results. We withdraw them as we
withdrew the first set.

What we are left with is a mechanism rather than a number. A provenance claim
that lives only in a docstring, a generated directory that becomes indistinguishable
from upstream material once it is inside the checkout, and a diversity statistic
that scores the forgery and the replacement alike: each is ordinary, and together
they were enough to make a team that had just written about fabricated data
fabricate it again. The checks in Section~\ref{sec:whytwice} are the ones we now
apply, and we would rather publish them with this history attached than publish
the benchmark result we thought we had.

%% ===========================================================================
\appendix

\section{Reproduction}
\label{app:repro}

Every number in this paper is produced by running the scripts below against the
result files below. Nothing is transcribed by hand, and the check at the end of
this appendix is the one we use to verify that claim.

\subsection{Inputs}

Table~\ref{tab:inputs} lists them.

\begin{table}[h]
\centering
\caption{Result files this paper reads. Sizes are of the released artefacts.}
\label{tab:inputs}
\footnotesize
\setlength{\tabcolsep}{4pt}
\begin{tabular}{@{}>{\raggedright\arraybackslash}p{0.50\linewidth}>{\raggedright\arraybackslash}p{0.42\linewidth}@{}}
\toprule
Path under \texttt{results/} & Contents \\
\midrule
\texttt{pilot/alerts.csv} &
  the \RNtp{} true-positive alerts: rule, message, file, line \\
\texttt{real\_fp/real\_fp\_alerts.json} &
  the \RNfp{} false-positive alerts, with the source corpus of each \\
\texttt{real\_fp/secure\_fp.sarif} &
  the raw CodeQL SARIF output the false positives were extracted from \\
\texttt{real\_eval/real\_eval\_results.json} &
  per-triager confusion counts, bootstrap MCC intervals, burden, and the
  \(12\) pairwise McNemar tests \\
\texttt{audit/audit\_results.json} &
  the \Nrubric{} explanation-rubric items and \Nprobe{} contamination probes \\
\bottomrule
\end{tabular}
\end{table}

\subsection{Pipeline}

Three stages, in order. The first consumes the corpus and produces the alert
inventory; the second consumes the inventory and produces the metrics; the third
consumes the metrics and produces this document.

\begin{enumerate}
\item \texttt{fix1\_real\_fp\_scan.py} runs the
\texttt{python-security-extended} query suite over two directories inside our
\texttt{SecurityEval} checkout, \texttt{Testcases\_Secure\_FP} and
\texttt{Testcases\_Secure\_Patched}, and writes every alert to
\texttt{real\_fp/real\_fp\_alerts.json} with the SARIF it came from. Its
docstring asserts that the benchmark's authors certify both directories as
secure. They do not exist upstream, and nobody certifies them. The script is
released unmodified, because the gap between what it claims and what it does is
the subject of this paper.

\item \texttt{fix2\_rebuild\_and\_triage.py} builds the two evaluation sets,
runs the four triagers, and writes
\texttt{real\_eval/real\_eval\_results.json}. It fixes the random seed to
\(42\); both the balanced set's choice of which false positives to keep and the
cross-validation folds derive from that seed.

\item \texttt{make\_tables.py} then \texttt{build.py} read the result files and
emit every macro, table row and figure used here. The composition and
clause-attribution analyses of Sections~\ref{sec:inventory}
and~\ref{sec:heurwhy} are computed in this stage directly from the two alert
files, and involve no sampling, so they are exact rather than seed-dependent.
\end{enumerate}

\subsection{What is and is not seed-dependent}

The distinction matters for anyone re-running this. The confusion counts,
bootstrap intervals, burden figures and McNemar tests of
Section~\ref{sec:realresults} all depend on the seed, through the balanced
set's subsample and the cross-validation split. The inventory composition of
Section~\ref{sec:inventory} and the clause attribution of
Section~\ref{sec:heurwhy} depend on neither: they are functions of the two
alert files alone and will reproduce bit-for-bit from any environment that can
read JSON and CSV.

\subsection{Verification}

We check the substitution rather than trust it. \texttt{make\_tables.py} writes
every macro from a computed value, and \texttt{build.py} refuses to emit
\texttt{main.tex} if a \texttt{\%\%MACROS\%\%} or \texttt{\%\%ROWS:\%\%} marker
has no corresponding generated file. A number that appears in the text without a
macro behind it is therefore a number we typed, and the only such numbers in
this paper are the two structural constants stated in prose: the \(5\) folds of
the cross-validation and the \(50\times\) multiplier of the burden metric, both
of which are choices rather than measurements.

\section*{CRediT authorship contribution statement}

\textbf{Vinh Thinh Le:} Conceptualization, Methodology, Writing -- original draft.
\textbf{Son X. Ha:} Conceptualization, Formal analysis, Investigation, Writing -- original draft, Writing -- review \& editing.
\textbf{Nguyen Ngoc Phien:} Supervision, Project administration, Writing -- review \& editing.
\textbf{Tung Q. Nguyen:} Software, Data curation, Visualization, Investigation.

\section*{Declaration of competing interest}

The authors declare that they have no known competing financial interests or
personal relationships that could have appeared to influence the work reported
in this paper.

\section*{Funding}

This research received no specific grant from any funding agency in the
public, commercial or not-for-profit sectors.

\section*{Data availability}

The measured result files, the generator scripts that read them, and the LaTeX
sources that substitute their output into this document are released with the
paper. Every number, table and figure here is produced by running those scripts
against those files; none is transcribed by hand.

The artefact (result files, generator scripts and \LaTeX{} sources) is not
publicly posted at the time of submission. We will deposit it in a public
repository under a persistent identifier upon acceptance, and we can supply it
to the editorial office at any point during review.

\bibliographystyle{elsarticle-harv}
\bibliography{refs}

\end{document}
